\documentclass[12pt,a4paper]{article}

% ================= PACKAGES =================

\usepackage[a4paper,margin=1in]{geometry}
\usepackage{amsmath,amssymb}
\usepackage{graphicx}
\usepackage{booktabs}
\usepackage{array}
\usepackage{longtable}
\usepackage{multirow}
\usepackage{titlesec}
\usepackage{xcolor}
\usepackage{hyperref}
\usepackage{setspace}
\usepackage{fancyhdr}

% ================= FONT =================

\usepackage{newtxtext,newtxmath}

% ================= LINE SPACING =================

\onehalfspacing

% ================= HEADER =================

\pagestyle{fancy}
\fancyhf{}
\fancyhead[L]{AI-Based Traffic Prediction}
\fancyhead[R]{Research Paper Template}
\fancyfoot[C]{\thepage}

% ================= SECTION STYLE =================

\titleformat{\section}
{\Large\bfseries\color{blue}}
{\thesection}{1em}{}

% ================= DOCUMENT =================

\begin{document}
	
	% ================= TITLE PAGE =================
	
	\begin{titlepage}
		
		\centering
		
		\vspace*{2cm}
		
		{\Huge\bfseries
			AI-Based Smart Traffic Congestion Prediction
			\par}
		
		\vspace{1cm}
		
		{\Large
			Using Deep Learning and IoT Systems
			\par}
		
		\vspace{2cm}
		
		\includegraphics[width=0.25\textwidth]{example-image}
		
		\vfill
		
		{\Large
			Brintha Theboral
			\par}
		
		\vspace{0.5cm}
		
		Department of Computer Science Engineering
		
		ABC Institute of Technology
		
		Chennai, India
		
		\vfill
		
		{\large \today}
		
	\end{titlepage}
	
	% ================= ABSTRACT =================
	
	\section*{Abstract}
	
	Urban traffic congestion has become a major issue in smart cities. This paper proposes an AI-driven traffic prediction framework using deep learning and IoT sensor systems. The proposed model improves traffic forecasting accuracy and enables efficient congestion management.
	
	\vspace{0.5cm}
	
	\textbf{Keywords:}
	Artificial Intelligence, Traffic Prediction, Smart Cities, Deep Learning, IoT
	
	% ================= INTRODUCTION =================
	
	\section{Introduction}
	
	Rapid urbanization has significantly increased traffic congestion worldwide. Intelligent transportation systems are essential for improving urban mobility and reducing fuel consumption.
	
	Deep learning techniques provide improved prediction capabilities compared to traditional statistical approaches.
	
	% ================= LITERATURE SURVEY =================
	
	\section{Literature Survey}
	
	Several studies have explored machine learning methods for traffic prediction.
	
	\begin{itemize}
		\item CNN-based traffic forecasting systems
		\item IoT-enabled smart transportation
		\item LSTM recurrent neural networks
		\item Hybrid AI optimization models
	\end{itemize}
	
	% ================= PROPOSED SYSTEM =================
	
	\section{Proposed System}
	
	The proposed system consists of:
	
	\begin{enumerate}
		\item Traffic Data Collection
		\item Sensor Integration
		\item Data Preprocessing
		\item AI Prediction Engine
		\item Real-Time Monitoring Dashboard
	\end{enumerate}
	
	% ================= EQUATIONS =================
	
	\section{Mathematical Model}
	
	Traffic flow prediction can be represented as:
	
	\[
	Y_t = f(X_t,W,b)
	\]
	
	where:
	
	\begin{itemize}
		\item $X_t$ = Input traffic data
		\item $W$ = Weight parameters
		\item $b$ = Bias
	\end{itemize}
	
	Loss function:
	
	\[
	L = \frac{1}{n}\sum_{i=1}^{n}(y_i - \hat{y_i})^2
	\]
	
	% ================= TABLE =================
	
	\section{Experimental Results}
	
	\begin{table}[h]
		\centering
		
		\caption{Performance Comparison}
		
		\begin{tabular}{|c|c|c|}
			\hline
			Model & Accuracy & RMSE \\
			\hline
			Linear Regression & 78.2\% & 12.4 \\
			\hline
			Random Forest & 85.6\% & 8.9 \\
			\hline
			LSTM & 93.1\% & 4.3 \\
			\hline
			CNN-LSTM Hybrid & 95.4\% & 3.1 \\
			\hline
		\end{tabular}
		
	\end{table}
	
	% ================= FIGURE =================
	
	\section{System Architecture}
	
	\begin{figure}[h]
		
		\centering
		
		\includegraphics[width=0.8\textwidth]{example-image}
		
		\caption{AI-Based Traffic Prediction Architecture}
		
	\end{figure}
	
	% ================= DISCUSSION =================
	
	\section{Discussion}
	
	The proposed hybrid deep learning model achieved higher prediction accuracy due to temporal feature extraction and real-time sensor integration.
	
	% ================= FUTURE WORK =================
	
	\section{Future Scope}
	
	Future enhancements include:
	
	\begin{itemize}
		\item Edge AI deployment
		\item 5G-enabled traffic systems
		\item Autonomous vehicle integration
		\item Smart signal optimization
	\end{itemize}
	
	% ================= CONCLUSION =================
	
	\section{Conclusion}
	
	This research demonstrated the effectiveness of AI-based traffic prediction systems for smart city infrastructure. The proposed framework improves scalability and prediction accuracy.
	
	% ================= REFERENCES =================
	
	\section*{References}
	
	\begin{enumerate}
		
		\item Smith J., ``Smart Traffic Monitoring Using IoT'', IEEE Transactions, 2023.
		
		\item Johnson A., ``Deep Learning for Traffic Forecasting'', Springer, 2022.
		
		\item Kumar R., ``Urban AI Systems'', Elsevier, 2021.
		
	\end{enumerate}
	
	% ================= APPENDIX =================
	
	\appendix
	
	\section{Dataset Details}
	
	The dataset contains:
	\begin{itemize}
		\item Vehicle count
		\item Signal timing
		\item GPS movement
		\item Weather conditions
	\end{itemize}
	
\end{document}
